\section{陈与隋：在长江与关中之间重新接合天下}
\label{sec:chen-sui-reunification}

五五七年陈霸先受禅建陈，首先接收的是梁末的州郡文书、军人、侨民和受损的建康周边秩序。陈文帝面对王琳等长江中游势力及北齐的介入，须依靠江州、湘州、京口和建康的军镇维持防线。五六二年王琳集团失败，部众和随军人口北迁，江汉、淮南的流民与军户分布随之改变。政权在建康恢复，并不等于长江流域的人、粮和道路恢复到侯景之乱以前的状态。

\subsection{陈朝的有限经营}

五六九年陈顼即位为宣帝，五七三年陈朝北伐，吴明彻等取淮南、寿阳附近城镇，并在新取地区安排守将、户籍和粮运。城镇一度易手可以由《陈书》《南史》及《资治通鉴》校核，实际控制能延伸到哪些乡里、地方对新任命如何反应，则不能仅由战报推断。陈的北伐表明建康仍有能力越过长江组织水陆军，也暴露新占地区必须立刻面对供给与安置的限制。

五八二年宣帝死，陈叔宝即位。后世史家对陈后主的奢侈、宫廷文学和近侍政治有强烈批评，不能把这些道德化画像直接当成陈亡的完整解释。长江下游仍维持水军、城戍和漕运，并承担持续军费；宫廷选择、地方财政和边防压力是同时存在的条件。陈并非没有国家能力，而是难以在北方重新整合后持续扩大其可用资源。

\subsection{隋在北方重造中枢}

\eventanchor{NSL-440-589-271}{五八一年杨坚受禅建隋}，重整中央官制、府兵与州郡任命，并削弱北周宗室影响。五八二年营建大兴城，宫城、官署、坊市、水源和道路的规划需要吸纳关中工匠与劳力；城市并非单靠图纸出现，而要靠运输、供给与日常劳作持续形成。面对突厥南下，隋以长城、州镇、机动军队和利用可汗间竞争的外交同时应对，说明北方统一仍受边境关系牵制。

五八三年开皇律令的修订、州郡整并、户籍清理与租调、均田框架的推进，五八四年广通渠连接大兴、潼关和黄河运输，都是把法律、赋役、道路和仓储接到同一中枢的尝试。制度条文与工程存在可以核实，实际实施的地域、户等负担和地方社会的反应却不能由《隋书》一类国家记述直接推出。五八七年废后梁、安置江陵官员与军人，进一步使隋取得荆襄道路和粮运资源，也使南下不再只是跨越长江的一次军事冒险。

\subsection{渡江之后，统一仍须安置}

五八八年隋以杨广、杨俊、杨素等分路伐陈，陈朝调集水军和沿岸守将；\eventanchor{NSL-440-589-290}{次年隋军入建康，陈后主被俘，陈亡}。统帅、出兵与降附的主线可靠，兵数、路线、长江渡口的每一细节和战后民众态度不宜以正史数字写实。隋随即接收州郡、户籍、仓储、官员、陈室宗族与江南士族，派使安抚并调整封授、迁置。军事胜利解决了最高政权的归属，却没有让江南社会、赋役和地方网络立即同北方完全一致。

约在这一时代，颜之推完成或修订《颜氏家训》的主要篇章，以北齐、北周、隋之间的家族经历反思读书、婚姻、仕宦和迁徙。成书、抄写与所述经验的时间并不必然重合，文本也不能代表所有家庭；它仍提醒我们，再统一并非只有皇帝、城池和律令，也是一代人在改朝换代中重新安排家计、教育与身份的过程。

\subsection{材料的边界}

陈末以《陈书》《南史》为主，隋初以《隋书》纪志及《北史》为主，并与《资治通鉴》交叉；这些材料都以最终统一为已知结果书写早期选择。城址、墓志、运河和仓储遗存，以及《颜氏家训》这样的家族文本，分别提供城市、官职、运输和生活经验的有限窗口。它们支持我们追问统一如何经由户籍、粮道、寺院、军队和家庭来实现，同时也要求避免将法律、工程或一位皇帝的决断误作整个社会已被无缝整合的证据。


% 年序补线：本节将既有账本中尚未初次落入正文的条目接入相应历史场景。
\subsection{新朝门槛：北齐、北周与陈的接收难题}
557年，陈围绕\eventanchor{NS-LEDGER-025-1}{陈霸先建陈}、\eventanchor{NS-LEDGER-025-2}{陈霸先建陈}、\eventanchor{NS-LEDGER-025-3}{陈霸先建陈}、\eventanchor{NS-LEDGER-025-4}{陈霸先建陈}、\eventanchor{NS-LEDGER-025-5}{陈霸先建陈}、\eventanchor{NS-LEDGER-025-6}{“陈霸先建陈”的异源材料与影响范围的复核}留下的其实是同一事件簇：陈霸先建陈。 现有材料只能把这一变化放在连续年序中；前期条件、命令响应、地方执行、后续调整和异文问题是同一事件的不同读法，不能伪装成六场独立历史。 《宋书》《南齐书》《梁书》《陈书》帝纪及列传；《魏书》《北齐书》《周书》；《资治通鉴》宋齐梁陈纪提供该簇的最低证据链；废立、即位或政权转换主干可核；礼仪、动机和清洗范围常受胜者叙事影响。

557年，陈的\eventanchor{NSL-440-589-222}{陈霸先受禅建立陈，梁朝终结}使陈霸先受禅建立陈，梁朝终结成为可追索的行动。 行动者受陈霸先掌握建康军队和政令，梁敬帝失去实际支持限制，随后陈成为江南新政权，面对战后人口、财政和边防重建。 材料的可说范围由《陈书》帝纪及列传；《南史》；《资治通鉴》界定；本纪和列传可互校；政变动机与参与范围常受胜者叙事影响。

\eventanchor{NSL-440-589-223}{陈朝接收梁末州郡文书、军人和侨民，重整建康周边秩序}所标出的陈朝接收梁末州郡文书、军人和侨民，重整建康周边秩序，发生在557年的陈。 将其置回事件链，前提是侯景之乱和梁末内战使行政记录与户籍严重破损，而后续处置面对陈的有效统治首先集中于长江下游，地方恢复不均衡。 所据《陈书》帝纪及列传；《南史》；《资治通鉴》留下可核的线索，同时提醒读者地方活动见地理、墓志或文书线索；精英材料不代表整体社会。

558年，陈的\eventanchor{NSL-440-589-225}{陈文帝即位后继续压制王琳等长江中游势力}写下陈文帝即位后继续压制王琳等长江中游势力这一处具体变化。 这里的资源与选择关系可从陈霸先去世使新朝必须迅速安抚军将和地方州镇看出，结果使陈朝逐步守住江东核心区，但荆湘和淮南仍不稳定。 材料来自《陈书》帝纪及列传；《南史》；《资治通鉴》；因此战事年月可据编年材料核对；兵数、杀伤与路线须保留史书夸饰风险。

在559年这一个节点，陈通过\eventanchor{NSL-440-589-226}{王琳在长江中游拥立梁永嘉王萧庄并联络北齐}显出王琳在长江中游拥立梁永嘉王萧庄并联络北齐的压力。 前因和后效须分开读：前者是梁末宗室与地方将领不愿接受陈朝取代梁室，后者是陈的统一江南努力受北齐和地方军府牵制。 以《陈书》帝纪及列传；《南史》；《资治通鉴》交叉检验，能够确认的止于此；战事年月可据编年材料核对；兵数、杀伤与路线须保留史书夸饰风险。

从559年的记录看，陈经由\eventanchor{NSL-440-589-298}{陈朝在建康战后修复宫城、仓储和长江水运节点}处理陈朝在建康战后修复宫城、仓储和长江水运节点。 此处不把时间相邻当作因果；能追到的是侯景之乱和梁末战争破坏都城基础设施与行政记录，以及其后江东恢复依赖州郡转运和地方劳力，重建过程并不均衡。 从《陈书》帝纪及列传；《南史》；《资治通鉴》可稳妥写到这一节点；地方活动见地理、墓志或文书线索；精英材料不代表整体社会。

\subsection{北周与陈：军政重组和江南恢复}
从560年的记录看，陈经由\eventanchor{NSL-440-589-229}{陈文帝继续以江州、湘州等军镇防御王琳与北齐压力}处理陈文帝继续以江州、湘州等军镇防御王琳与北齐压力。 它不是孤立转折：长江中游仍有梁室残余和北齐援军构成背景，陈朝军费和军粮优先投向江防，地方行政恢复受限则把压力送往下一阶段。 就材料边界说，《陈书》帝纪及列传；《南史》；《资治通鉴》支持这一轮廓；战事年月可据编年材料核对；兵数、杀伤与路线须保留史书夸饰风险。

562年，陈的\eventanchor{NSL-440-589-232}{陈军击败王琳集团，萧庄与王琳势力北走}写下陈军击败王琳集团，萧庄与王琳势力北走这一处具体变化。 将其置回事件链，前提是陈朝集中江东、江南兵力，切断对手长江水路联络，而后续处置面对陈在长江中下游的统治得到巩固，北齐仍控制江北部分地区。 《陈书》帝纪及列传；《南史》；《资治通鉴》可互校时间位置与主体，但战事年月可据编年材料核对；兵数、杀伤与路线须保留史书夸饰风险。

562年的年序到\eventanchor{NSL-440-589-233}{王琳部众与随军人口北迁，改变江汉、淮南的流民与军户分布}时，陈须面对王琳部众与随军人口北迁，改变江汉、淮南的流民与军户分布。 这一行动的条件是地方军府失败后将领、家属和部曲需重新寻找保护者；其直接余波则是战败者的迁徙为北齐、陈两方提供或增加安置负担。 证据不允许把场景填满：《陈书》帝纪及列传；《南史》；《资治通鉴》仅能支持核心事项，地方活动见地理、墓志或文书线索；精英材料不代表整体社会。

在566年这一个节点，陈通过\eventanchor{NSL-440-589-240}{陈文帝病逝，废帝即位，朝廷由重臣辅政}显出陈文帝病逝，废帝即位，朝廷由重臣辅政的压力。 这一节点将皇位继承时新君年幼，陈朝需在边境压力下维持中枢稳定与陈顼等宗室和辅臣的政治力量上升接合，却不预设两者之间只有一种因果。 从《陈书》帝纪及列传；《南史》；《资治通鉴》可稳妥写到这一节点；本纪和列传可互校；政变动机与参与范围常受胜者叙事影响。

把\eventanchor{NSL-440-589-241}{陈顼控制朝政并处理建康宗室、军将与地方任命}放回567年，可见陈正围绕陈顼控制朝政并处理建康宗室、军将与地方任命作出处理。 其代价落在幼主政局为强宗室介入提供机会所限定的处境中，并以陈朝权力进一步集中于陈顼集团，皇统再度脆弱延续下去。 材料的可说范围由《陈书》帝纪及列传；《南史》；《资治通鉴》界定；本纪和列传可互校；政变动机与参与范围常受胜者叙事影响。

568年的年序到\eventanchor{NSL-440-589-242}{陈顼废陈伯宗为临海王，独揽朝政}时，陈须面对陈顼废陈伯宗为临海王，独揽朝政。 它承接陈朝内廷围绕幼主、宗室与辅臣冲突加剧，并使陈顼在受禅前完成对禁军与政令的控制。 以《陈书》帝纪及列传；《南史》；《资治通鉴》交叉检验，能够确认的止于此；本纪和列传可互校；政变动机与参与范围常受胜者叙事影响。

569年，陈的\eventanchor{NSL-440-589-243}{陈顼即皇帝位，是为宣帝}写下陈顼即皇帝位，是为宣帝这一处具体变化。 行动者受辅政者已控制禁军和诏令，幼主缺乏独立支持限制，随后陈朝在较强皇帝领导下恢复对江南州郡的整合。 《陈书》帝纪及列传；《南史》；《资治通鉴》保留了这条年序，然而本纪和列传可互校；政变动机与参与范围常受胜者叙事影响。

\subsection{灭齐与边境：北方统一尚待完成}
573年的材料将陈的\eventanchor{NSL-440-589-249}{陈宣帝北伐，吴明彻等攻取淮南、寿阳附近多处城镇}系于陈宣帝北伐，吴明彻等攻取淮南、寿阳附近多处城镇，其影响由此进入后续年序。 从中可辨认出北齐内廷紊乱且江北防线薄弱，陈欲恢复淮南资源的约束，及其对陈一度扩张至淮河，需面对远距离驻军和粮运困难的影响。 《陈书》帝纪及列传；《南史》；《资治通鉴》可互校时间位置与主体，但战事年月可据编年材料核对；兵数、杀伤与路线须保留史书夸饰风险。

573年，陈的\eventanchor{NSL-440-589-250}{陈朝在新取淮南地区安排守将、户籍和粮运}写下陈朝在新取淮南地区安排守将、户籍和粮运这一处具体变化。 这一行动的条件是军事占领若无州县与转运支持便难持续；其直接余波则是陈的北伐成果受地方整合能力制约，边防支出上升。 年序和核心行动以《陈书》帝纪及列传；《南史》；《资治通鉴》为据；制度条文可核；实际执行地域、户等与负担不可由诏令直接推定。

\subsection{隋陈之间：战争、接收与再统一的限度}
把\eventanchor{NSL-440-589-304}{尉迟迥之乱中，杨坚军利用关中与河北粮道协调平叛}放回580年，可见隋前夜正围绕尉迟迥之乱中，杨坚军利用关中与河北粮道协调平叛作出处理。 这一节点将叛乱跨越原北齐旧地，单靠禁军不足以完成镇压与平叛检验了新权力集团的后勤能力，并强化其对河北的控制接合，却不预设两者之间只有一种因果。 就材料边界说，《周书》帝纪及列传；《北史·周本纪》；《资治通鉴》支持这一轮廓；战事年月可据编年材料核对；兵数、杀伤与路线须保留史书夸饰风险。

把\eventanchor{NS-LEDGER-033-1}{杨坚受禅建隋}、\eventanchor{NS-LEDGER-033-2}{杨坚受禅建隋}、\eventanchor{NS-LEDGER-033-3}{杨坚受禅建隋}、\eventanchor{NS-LEDGER-033-4}{杨坚受禅建隋}、\eventanchor{NS-LEDGER-033-5}{杨坚受禅建隋}、\eventanchor{NS-LEDGER-033-6}{“杨坚受禅建隋”的异源材料与影响范围的复核}放在一起读，才能看见581年隋所经历的杨坚受禅建隋是一条事件链。 现有材料只能把这一变化放在连续年序中；前期条件、命令响应、地方执行、后续调整和异文问题是同一事件的不同读法，不能伪装成六场独立历史。 《宋书》《南齐书》《梁书》《陈书》帝纪及列传；《魏书》《北齐书》《周书》；《资治通鉴》宋齐梁陈纪提供该簇的最低证据链；废立、即位或政权转换主干可核；礼仪、动机和清洗范围常受胜者叙事影响。

581年，\eventanchor{NSL-440-589-272}{隋初重整中央官制、府兵与州郡任命，削减北周宗室影响}并非抽象名称，而是隋发生隋初重整中央官制、府兵与州郡任命，削减北周宗室影响的叙事入口。 它使原有的新朝需防止北周复辟并使军政体系服从新皇权转化为需要继续承担的隋的制度整合为再统一提供行政骨架，旧制仍有过渡痕迹。 《隋书·高祖纪》及相关志；《北史》；《资治通鉴》构成最低证据链，故制度条文可核；实际执行地域、户等与负担不可由诏令直接推定。

从城邑、军镇或官署的实际处境出发，581年的\eventanchor{NSL-440-589-273}{突厥沙钵略可汗支持北周宗室反隋，北部边境紧张}关乎隋与突厥的突厥沙钵略可汗支持北周宗室反隋，北部边境紧张。 此处不把时间相邻当作因果；能追到的是突厥担心隋统一北方后削弱其外交筹码，亦利用北周遗臣，以及其后隋须同时处置北方威胁和南陈关系，外交分化策略更受重视。 这里以《周书·突厥传》；《北史·突厥传》；《资治通鉴》为证据支点，且明确保留政权互动可据多方传记校核；族称、疆界和战果须避免按后世分类固化。

把\eventanchor{NSL-440-589-274}{陈宣帝去世，陈叔宝即位，是为后主}放回582年，可见陈正围绕陈宣帝去世，陈叔宝即位，是为后主作出处理。 它把陈朝继承在北方隋朝新建与边防压力下发生带入新的局面，令后主朝内廷政治增强，陈的军事与财政准备受到影响成为后来人的问题。 从《陈书》帝纪及列传；《南史》；《资治通鉴》可稳妥写到这一节点；本纪和列传可互校；政变动机与参与范围常受胜者叙事影响。

在582年这一个节点，隋通过\eventanchor{NSL-440-589-275}{隋营建大兴城，规划宫城、官署、坊市与交通供给}显出隋营建大兴城，规划宫城、官署、坊市与交通供给的压力。 前因和后效须分开读：前者是关中旧城条件不足以容纳新朝中枢和人口集聚，后者是新都建设集中工匠、木材和转运资源，成为隋唐长安城市基础。 所据《隋书·高祖纪》及相关志；《北史》；《资治通鉴》留下可核的线索，同时提醒读者地方活动见地理、墓志或文书线索；精英材料不代表整体社会。

582年，隋与突厥的\eventanchor{NSL-440-589-276}{突厥大举南下，隋以长城、州镇和机动军队应对}写下突厥大举南下，隋以长城、州镇和机动军队应对这一处具体变化。 将其置回事件链，前提是北周隋代之际的外交破裂使草原军事压力集中爆发，而后续处置面对隋在北方防务中学习分化可汗与重建边镇供给。 以《周书·突厥传》；《北史·突厥传》；《资治通鉴》交叉检验，能够确认的止于此；战事年月可据编年材料核对；兵数、杀伤与路线须保留史书夸饰风险。

在582年这一个节点，隋通过\eventanchor{NSL-440-589-305}{大兴城建设同步规划水源、道路与坊市，吸纳关中工匠和劳力}显出大兴城建设同步规划水源、道路与坊市，吸纳关中工匠和劳力的压力。 这一处理回应新都不仅是宫廷工程，也要承担行政、市场与军队供给功能，同时改变了城市建设强化关中中心性，并为隋唐长安留下空间框架。 这项判断依托《隋书·高祖纪》及相关志；《北史》；《资治通鉴》，并受限于地方活动见地理、墓志或文书线索；精英材料不代表整体社会。

在583年，隋的\eventanchor{NS-LEDGER-034-1}{隋整顿北方军政与户籍}、\eventanchor{NS-LEDGER-034-2}{隋整顿北方军政与户籍}、\eventanchor{NS-LEDGER-034-3}{隋整顿北方军政与户籍}、\eventanchor{NS-LEDGER-034-4}{隋整顿北方军政与户籍}、\eventanchor{NS-LEDGER-034-5}{隋整顿北方军政与户籍}、\eventanchor{NS-LEDGER-034-6}{“隋整顿北方军政与户籍”的异源材料与影响范围的复核}并非六件孤立大事，而是隋整顿北方军政与户籍的六个核验入口。 现有材料只能把这一变化放在连续年序中；前期条件、命令响应、地方执行、后续调整和异文问题是同一事件的不同读法，不能伪装成六场独立历史。 《宋书》《南齐书》《梁书》《陈书》帝纪及列传；《魏书》《北齐书》《周书》；《资治通鉴》宋齐梁陈纪提供该簇的最低证据链；制度或政令形成可追索；执行地域、对象和持续时间不能由条文直接推定。

583年的年序到\eventanchor{NSL-440-589-277}{隋修订开皇律令，强调简约刑法与中央审断}时，隋须面对隋修订开皇律令，强调简约刑法与中央审断。 这一行动的条件是新朝需要以统一法令取代北周末年的政令混乱；其直接余波则是法律规范成为整合新附州郡的工具，地方司法实践仍须由个案检验。 材料的可说范围由《隋书·高祖纪》及相关志；《北史》；《资治通鉴》界定；制度条文可核；实际执行地域、户等与负担不可由诏令直接推定。

\eventanchor{NSL-440-589-278}{隋整并州郡、清理户籍与赋役，推进均田和租调体系}把583年的隋与隋整并州郡、清理户籍与赋役，推进均田和租调体系这一材料所见的处置连在一起。 这一行动的条件是北周、北齐旧地制度不一，新朝需建立可征收的人户框架；其直接余波则是中央财政能力提高，但隐户、流民和地区差异不可能立即消失。 以《隋书·高祖纪》及相关志；《北史》；《资治通鉴》交叉检验，能够确认的止于此；制度条文可核；实际执行地域、户等与负担不可由诏令直接推定。

在583年的多方行动者之间，\eventanchor{NSL-440-589-279}{隋利用突厥可汗间竞争，削弱沙钵略对北周遗臣的支持}标记隋与突厥面对隋利用突厥可汗间竞争，削弱沙钵略对北周遗臣的支持的选择。 由此能看到的是突厥汗国内部权力并非统一，隋可用使节、婚姻和物资分化与北方压力暂获缓解，隋得以更多筹划南方战争之间的条件关系，而非事后必然论。 材料的可说范围由《周书·突厥传》；《北史·突厥传》；《资治通鉴》界定；政权互动可据多方传记校核；族称、疆界和战果须避免按后世分类固化。

584年，隋的\eventanchor{NSL-440-589-280}{隋开通广通渠，连接大兴城与潼关、黄河运输体系}使隋开通广通渠，连接大兴城与潼关、黄河运输体系成为可追索的行动。 它不是孤立转折：新都和军队需要稳定粮食、木材与人员水运构成背景，关中与东部仓储联系加强，为大规模南征提供后勤条件则把压力送往下一阶段。 年序和核心行动以《隋书·高祖纪》及相关志；《北史》；《资治通鉴》为据；制度条文可核；实际执行地域、户等与负担不可由诏令直接推定。

\eventanchor{NSL-440-589-281}{陈朝面对隋北方整合，修整长江防线并调整江北守将}所标出的陈朝面对隋北方整合，修整长江防线并调整江北守将，发生在584年的陈。 此处不把时间相邻当作因果；能追到的是隋统一北方后可集中资源南下，陈需防范渡江与淮河攻势，以及其后陈军政更依赖长江水道和少数重镇，战略纵深有限。 以《陈书》帝纪及列传；《南史》；《资治通鉴》交叉检验，能够确认的止于此；战事年月可据编年材料核对；兵数、杀伤与路线须保留史书夸饰风险。

从城邑、军镇或官署的实际处境出发，585年的\eventanchor{NSL-440-589-282}{隋在北方旧齐地区继续编户、核田与恢复仓储转运}关乎隋的隋在北方旧齐地区继续编户、核田与恢复仓储转运。 它使原有的州郡整并后的行政命令需落实到乡里和粮仓转化为需要继续承担的新附人口被纳入国家记录，基层社会的适应与规避并存。 证据不允许把场景填满：《隋书·高祖纪》及相关志；《北史》；《资治通鉴》仅能支持核心事项，地方活动见地理、墓志或文书线索；精英材料不代表整体社会。

585年，\eventanchor{NSL-440-589-283}{陈后主朝廷的奢侈、宫廷文学与近侍政治受到史家强烈批评}并非抽象名称，而是陈发生陈后主朝廷的奢侈、宫廷文学与近侍政治受到史家强烈批评的叙事入口。 前因和后效须分开读：前者是陈末宫廷文化与政治决策相互交织，后出亡国叙事常道德化，后者是应与财政、军镇和地方控制材料并读，不能以人物品评替代结构解释。 《陈书》帝纪及列传；《南史》；《资治通鉴》可互校时间位置与主体，但成书、抄写与所述事态须分开；传本年代和作者归属尚有可讨论处。

在586年这一个节点，隋通过\eventanchor{NSL-440-589-284}{隋继续整顿府兵、仓储和江淮方向的军情侦察}显出隋继续整顿府兵、仓储和江淮方向的军情侦察的压力。 它把南陈未灭且长江天险要求长期筹备水陆协同带入新的局面，令隋军逐步形成多路南下的后勤和指挥条件成为后来人的问题。 《隋书·高祖纪》及相关志；《北史》；《资治通鉴》使这项处置可被追查，不足以抹平战事年月可据编年材料核对；兵数、杀伤与路线须保留史书夸饰风险。

把\eventanchor{NSL-440-589-285}{陈朝在长江下游维持水军、城戍与漕运，承受持续军费}放回586年，可见陈正围绕陈朝在长江下游维持水军、城戍与漕运，承受持续军费作出处理。 这一行动的条件是隋的威胁要求陈将有限资源集中于江防；其直接余波则是财政压力与地方动员能力不足限制陈朝长期抵抗。 这项判断依托《陈书》帝纪及列传；《南史》；《资治通鉴》，并受限于制度条文可核；实际执行地域、户等与负担不可由诏令直接推定。

587年的年序到\eventanchor{NSL-440-589-286}{隋废后梁，接收荆襄州郡与江陵资源}时，隋须面对隋废后梁，接收荆襄州郡与江陵资源。 将其置回事件链，前提是后梁长期依附北周、隋，已难在隋南征计划中保持自主，而后续处置面对隋取得长江中游战略支点，对陈形成西北两面压力。 就材料边界说，《隋书·高祖纪》及相关志；《北史》；《资治通鉴》支持这一轮廓；本纪和列传可互校；政变动机与参与范围常受胜者叙事影响。

587年，隋的\eventanchor{NSL-440-589-287}{隋在荆襄安置后梁官员、军人和居民，整顿道路与粮运}写下隋在荆襄安置后梁官员、军人和居民，整顿道路与粮运这一处具体变化。 可见的选择边界来自新附地区必须迅速提供南征通道而又要避免地方反弹，并留下荆州成为隋军沿江推进的重要后方，整合成本由地方承担。 这里以《隋书·高祖纪》及相关志；《北史》；《资治通鉴》为证据支点，且明确保留地方活动见地理、墓志或文书线索；精英材料不代表整体社会。

对于588年的隋与陈而言，\eventanchor{NS-LEDGER-035-1}{隋军南下伐陈}、\eventanchor{NS-LEDGER-035-2}{隋军南下伐陈}、\eventanchor{NS-LEDGER-035-3}{隋军南下伐陈}、\eventanchor{NS-LEDGER-035-4}{隋军南下伐陈}、\eventanchor{NS-LEDGER-035-5}{隋军南下伐陈}、\eventanchor{NS-LEDGER-035-6}{“隋军南下伐陈”的异源材料与影响范围的复核}都回到同一具体节点：隋军南下伐陈。 现有材料只能把这一变化放在连续年序中；前期条件、命令响应、地方执行、后续调整和异文问题是同一事件的不同读法，不能伪装成六场独立历史。 《宋书》《南齐书》《梁书》《陈书》帝纪及列传；《魏书》《北齐书》《周书》；《资治通鉴》宋齐梁陈纪提供该簇的最低证据链；核心战事与大致年序可核；军数、战果、路线和地方承受须以异政权记录及考古材料复核。

在588年的多方行动者之间，\eventanchor{NSL-440-589-288}{隋以杨广、杨俊、杨素等分路部署，发动灭陈战争}标记隋面对隋以杨广、杨俊、杨素等分路部署，发动灭陈战争的选择。 这一行动的条件是北方、荆襄和淮南准备完成，隋判断陈朝难以协调防御；其直接余波则是隋军自多方向逼近长江，陈的守备被迫分散。 就材料边界说，《隋书·高祖纪》及相关志；《北史》；《资治通鉴》支持这一轮廓；战事年月可据编年材料核对；兵数、杀伤与路线须保留史书夸饰风险。

\eventanchor{NSL-440-589-289}{陈朝调集长江水军和沿岸守将抵御隋军}把588年的陈与陈朝调集长江水军和沿岸守将抵御隋军这一材料所见的处置连在一起。 这里的资源与选择关系可从隋多路南下威胁京口、采石与江陵方向看出，结果使陈军指挥与城防虽有抵抗，难以弥补兵源和后勤差距。 这里以《陈书》帝纪及列传；《南史》；《资治通鉴》为证据支点，且明确保留战事年月可据编年材料核对；兵数、杀伤与路线须保留史书夸饰风险。

589年的隋与陈，通过\eventanchor{NS-LEDGER-036-1}{建康失陷，陈亡}、\eventanchor{NS-LEDGER-036-2}{建康失陷，陈亡}、\eventanchor{NS-LEDGER-036-3}{建康失陷，陈亡}、\eventanchor{NS-LEDGER-036-4}{建康失陷，陈亡}、\eventanchor{NS-LEDGER-036-5}{建康失陷，陈亡}、\eventanchor{NS-LEDGER-036-6}{“建康失陷，陈亡”的异源材料与影响范围的复核}把建康失陷，陈亡这一历史过程从条件、处置到余波连成一体。 现有材料只能把这一变化放在连续年序中；前期条件、命令响应、地方执行、后续调整和异文问题是同一事件的不同读法，不能伪装成六场独立历史。 《宋书》《南齐书》《梁书》《陈书》帝纪及列传；《魏书》《北齐书》《周书》；《资治通鉴》宋齐梁陈纪提供该簇的最低证据链；核心战事与大致年序可核；军数、战果、路线和地方承受须以异政权记录及考古材料复核。

从城邑、军镇或官署的实际处境出发，589年的\eventanchor{NSL-440-589-291}{隋接收陈朝州郡、户籍、仓储和官员，派使安抚江南}关乎隋的隋接收陈朝州郡、户籍、仓储和官员，派使安抚江南。 这一处理回应军事征服若无行政接收便无法将江南资源纳入新朝，同时改变了隋初在江南保留和调整旧有地方网络，统一不等于立即同质化。 《隋书·高祖纪》及相关志；《北史》；《资治通鉴》可互校时间位置与主体，但地方活动见地理、墓志或文书线索；精英材料不代表整体社会。

\eventanchor{NSL-440-589-292}{隋处理陈室宗族、降将与江南士族的封授和迁置}所标出的隋处理陈室宗族、降将与江南士族的封授和迁置，发生在589年的隋。 可见的选择边界来自新朝需防范复辟并获取熟悉江南文书、道路与社会的合作人，并留下政治安置塑造隋初江南精英与中央的关系。 证据不允许把场景填满：《隋书·高祖纪》及相关志；《北史》；《资治通鉴》仅能支持核心事项，本纪和列传可互校；政变动机与参与范围常受胜者叙事影响。

约589年，北齐遗民与隋的\eventanchor{NSL-440-589-293}{颜之推晚年完成或修订《颜氏家训》的主要篇章，反思北齐、北周、隋之间的家族处境}使颜之推晚年完成或修订《颜氏家训》的主要篇章，反思北齐、北周、隋之间的家族处境成为可追索的行动。 这里的资源与选择关系可从长期战乱与多次政权转换促使士族以家训整理生存经验看出，结果使文本为观察北方士族教育、婚姻和迁徙提供线索，篇章年代不可一概而论。 以《颜氏家训》及序跋；《隋书·经籍志》；相关版本研究交叉检验，能够确认的止于此；成书、抄写与所述事态须分开；传本年代和作者归属尚有可讨论处。



\subsection{577—586年：材料所见的连续压力}
北周末至隋初，制度接收与陈朝经营并行；淮河、长江、岭南和北方道路的户籍、军粮与网络决定再统一如何筹划。

公元577年（年序桥接），陈、北齐与北周的\eventanchor{JH-DENS-577-A}{陈朝立足江南，北齐和北周竞争关东与关中；边防、河运、府兵与地方豪强的组织决定战争可持续性}把陈朝立足江南，北齐和北周竞争关东与关中；边防、河运、府兵与地方豪强的组织决定战争可持续性留在年序中；\eventanchor{JH-DENS-577-B}{陈、北齐与北周的流民、侨置户、军镇居民或地方家族的安置与编户问题构成社会背景}则从陈、北齐与北周的流民、侨置户、军镇居民或地方家族的安置与编户问题构成社会背景观察同一年的约束。 这两项观察不宣称另有一场未载战争或政策，只把已有压力的时间位置和可见面向分开。 两种材料不能互相替代：《陈书》武文宣帝纪；《北齐书》《周书》；《资治通鉴》陈纪所能确认的有限，本纪与编年材料主要确证政权年序和精英行动；对基层执行、疆界和人口数量的沉默不得以叙事补足；《周书》武帝纪；墓志、造像、城址与道路考古同样只照见一部分，墓志、家族文本和侨置名目各自有选择性，不能据此统计迁徙规模或概括整体社会态度。

公元578年（年序桥接），\eventanchor{JH-DENS-578-A}{北周末隋初的制度接收与陈朝江南经营并行，淮河、长江和岭南水陆通道、户籍与军粮为再统一提供条件}和\eventanchor{JH-DENS-578-B}{隋与陈及突厥、岭南地方集团的寺院、僧侣、道士和施主网络在材料中连接都城、州郡与边地}共同避免把隋与陈及突厥、岭南地方集团压缩为单一都城事件：前者关注北周末隋初的制度接收与陈朝江南经营并行，淮河、长江和岭南水陆通道、户籍与军粮为再统一提供条件，后者关注隋与陈及突厥、岭南地方集团的寺院、僧侣、道士和施主网络在材料中连接都城、州郡与边地。 这使读者能把下一年的军事、行政或社会变化放回尚未消散的限制中。 前者依据《隋书·高祖纪》；《陈书·后主纪》；《北史》；《资治通鉴》陈纪；本纪与编年材料主要确证政权年序和精英行动；对基层执行、疆界和人口数量的沉默不得以叙事补足。后者参照《隋书》食货地理二志；墓志、寺院题记与江南城址；僧传、道教文本、造像记与遗址的成书或营造年代须分开，个案不能量化信仰或政策落实。

把公元579年（年序桥接）放回连续叙事，\eventanchor{JH-DENS-579-A}{北周末隋初的制度接收与陈朝江南经营并行，淮河、长江和岭南水陆通道、户籍与军粮为再统一提供条件}保留北周末隋初的制度接收与陈朝江南经营并行，淮河、长江和岭南水陆通道、户籍与军粮为再统一提供条件，\eventanchor{JH-DENS-579-B}{隋与陈及突厥、岭南地方集团的关隘、渡口、军镇和水陆道路的可通行性继续约束使节、军队和物资的行动}让隋与陈及突厥、岭南地方集团的关隘、渡口、军镇和水陆道路的可通行性继续约束使节、军队和物资的行动继续约束隋与陈及突厥、岭南地方集团的行动。 如此处理不是用空白填满史实，而是避免后来的转折抹去本年仍在运作的条件。 此处的证据支点是《隋书·高祖纪》；《陈书·后主纪》；《北史》；《资治通鉴》陈纪和《隋书》食货地理二志；墓志、寺院题记与江南城址；本纪与编年材料主要确证政权年序和精英行动；对基层执行、疆界和人口数量的沉默不得以叙事补足。，同时纪传中的行军路线、兵数与控制范围常被简化或夸饰，地理材料只能提供有限交叉。

\eventanchor{JH-DENS-580-A}{北周末隋初的制度接收与陈朝江南经营并行，淮河、长江和岭南水陆通道、户籍与军粮为再统一提供条件}和\eventanchor{JH-DENS-580-B}{隋与陈及突厥、岭南地方集团的朝廷和州郡围绕户籍、田租、军粮与转运的安排进入材料}使公元580年（年序桥接）的隋与陈及突厥、岭南地方集团既保留北周末隋初的制度接收与陈朝江南经营并行，淮河、长江和岭南水陆通道、户籍与军粮为再统一提供条件，也不略去隋与陈及突厥、岭南地方集团的朝廷和州郡围绕户籍、田租、军粮与转运的安排进入材料。 如此处理不是用空白填满史实，而是避免后来的转折抹去本年仍在运作的条件。 就可说范围而言，前一判断止于《隋书·高祖纪》；《陈书·后主纪》；《北史》；《资治通鉴》陈纪，本纪与编年材料主要确证政权年序和精英行动；对基层执行、疆界和人口数量的沉默不得以叙事补足；后一判断止于《隋书》食货地理二志；墓志、寺院题记与江南城址，志书、诏令与本纪可显示制度设计或战争需求，不能直接换算赋额、仓储或实际负担。

在公元581年（年序桥接），\eventanchor{JH-DENS-581-A}{北周末隋初的制度接收与陈朝江南经营并行，淮河、长江和岭南水陆通道、户籍与军粮为再统一提供条件}所见的北周末隋初的制度接收与陈朝江南经营并行，淮河、长江和岭南水陆通道、户籍与军粮为再统一提供条件与\eventanchor{JH-DENS-581-B}{隋与陈及突厥、岭南地方集团的流民、侨置户、军镇居民或地方家族的安置与编户问题构成社会背景}所见的隋与陈及突厥、岭南地方集团的流民、侨置户、军镇居民或地方家族的安置与编户问题构成社会背景并行发生在隋与陈及突厥、岭南地方集团的历史空间。 这里保留的是可核的连续性，而非对基层反应、疆界或人数的无据补写。 证据不许把这两个观察写成全域事实：《隋书·高祖纪》；《陈书·后主纪》；《北史》；《资治通鉴》陈纪受本纪与编年材料主要确证政权年序和精英行动；对基层执行、疆界和人口数量的沉默不得以叙事补足。限制，《隋书》食货地理二志；墓志、寺院题记与江南城址亦受墓志、家族文本和侨置名目各自有选择性，不能据此统计迁徙规模或概括整体社会态度限制。

公元582年（年序桥接），隋与陈及突厥、岭南地方集团的\eventanchor{JH-DENS-582-A}{北周末隋初的制度接收与陈朝江南经营并行，淮河、长江和岭南水陆通道、户籍与军粮为再统一提供条件}把北周末隋初的制度接收与陈朝江南经营并行，淮河、长江和岭南水陆通道、户籍与军粮为再统一提供条件留在年序中；\eventanchor{JH-DENS-582-B}{隋与陈及突厥、岭南地方集团的寺院、僧侣、道士和施主网络在材料中连接都城、州郡与边地}则从隋与陈及突厥、岭南地方集团的寺院、僧侣、道士和施主网络在材料中连接都城、州郡与边地观察同一年的约束。 它们不是由年份自动推出的因果，而是让前后事件之间的资源、道路、户籍或地方网络保持可见。 《隋书·高祖纪》；《陈书·后主纪》；《北史》；《资治通鉴》陈纪只能确认前一观察的基本年序，本纪与编年材料主要确证政权年序和精英行动；对基层执行、疆界和人口数量的沉默不得以叙事补足；《隋书》食货地理二志；墓志、寺院题记与江南城址为后一观察提供交叉线索，僧传、道教文本、造像记与遗址的成书或营造年代须分开，个案不能量化信仰或政策落实。

\eventanchor{JH-DENS-583-A}{北周末隋初的制度接收与陈朝江南经营并行，淮河、长江和岭南水陆通道、户籍与军粮为再统一提供条件}在公元583年（年序桥接）提示北周末隋初的制度接收与陈朝江南经营并行，淮河、长江和岭南水陆通道、户籍与军粮为再统一提供条件仍未结案；\eventanchor{JH-DENS-583-B}{隋与陈及突厥、岭南地方集团的关隘、渡口、军镇和水陆道路的可通行性继续约束使节、军队和物资的行动}以隋与陈及突厥、岭南地方集团的关隘、渡口、军镇和水陆道路的可通行性继续约束使节、军队和物资的行动说明隋与陈及突厥、岭南地方集团的选择还受具体条件限制。 如此处理不是用空白填满史实，而是避免后来的转折抹去本年仍在运作的条件。 材料边界须分别保留：《隋书·高祖纪》；《陈书·后主纪》；《北史》；《资治通鉴》陈纪与《隋书》食货地理二志；墓志、寺院题记与江南城址所见不同，且本纪与编年材料主要确证政权年序和精英行动；对基层执行、疆界和人口数量的沉默不得以叙事补足；纪传中的行军路线、兵数与控制范围常被简化或夸饰，地理材料只能提供有限交叉。

把公元584年（年序桥接）放回连续叙事，\eventanchor{JH-DENS-584-A}{北周末隋初的制度接收与陈朝江南经营并行，淮河、长江和岭南水陆通道、户籍与军粮为再统一提供条件}保留北周末隋初的制度接收与陈朝江南经营并行，淮河、长江和岭南水陆通道、户籍与军粮为再统一提供条件，\eventanchor{JH-DENS-584-B}{隋与陈及突厥、岭南地方集团的朝廷和州郡围绕户籍、田租、军粮与转运的安排进入材料}让隋与陈及突厥、岭南地方集团的朝廷和州郡围绕户籍、田租、军粮与转运的安排进入材料继续约束隋与陈及突厥、岭南地方集团的行动。 这两项观察不宣称另有一场未载战争或政策，只把已有压力的时间位置和可见面向分开。 就可说范围而言，前一判断止于《隋书·高祖纪》；《陈书·后主纪》；《北史》；《资治通鉴》陈纪，本纪与编年材料主要确证政权年序和精英行动；对基层执行、疆界和人口数量的沉默不得以叙事补足；后一判断止于《隋书》食货地理二志；墓志、寺院题记与江南城址，志书、诏令与本纪可显示制度设计或战争需求，不能直接换算赋额、仓储或实际负担。

公元585年（年序桥接），\eventanchor{JH-DENS-585-A}{北周末隋初的制度接收与陈朝江南经营并行，淮河、长江和岭南水陆通道、户籍与军粮为再统一提供条件}和\eventanchor{JH-DENS-585-B}{隋与陈及突厥、岭南地方集团的流民、侨置户、军镇居民或地方家族的安置与编户问题构成社会背景}共同避免把隋与陈及突厥、岭南地方集团压缩为单一都城事件：前者关注北周末隋初的制度接收与陈朝江南经营并行，淮河、长江和岭南水陆通道、户籍与军粮为再统一提供条件，后者关注隋与陈及突厥、岭南地方集团的流民、侨置户、军镇居民或地方家族的安置与编户问题构成社会背景。 它们将政治时间同资源和地方网络并置，却不把并置误作已被证实的直接因果。 材料边界须分别保留：《隋书·高祖纪》；《陈书·后主纪》；《北史》；《资治通鉴》陈纪与《隋书》食货地理二志；墓志、寺院题记与江南城址所见不同，且本纪与编年材料主要确证政权年序和精英行动；对基层执行、疆界和人口数量的沉默不得以叙事补足；墓志、家族文本和侨置名目各自有选择性，不能据此统计迁徙规模或概括整体社会态度。

在公元586年（年序桥接）的并行行动者之间，\eventanchor{JH-DENS-586-A}{北周末隋初的制度接收与陈朝江南经营并行，淮河、长江和岭南水陆通道、户籍与军粮为再统一提供条件}将北周末隋初的制度接收与陈朝江南经营并行，淮河、长江和岭南水陆通道、户籍与军粮为再统一提供条件落到时间位置，\eventanchor{JH-DENS-586-B}{隋与陈及突厥、岭南地方集团的寺院、僧侣、道士和施主网络在材料中连接都城、州郡与边地}从隋与陈及突厥、岭南地方集团的寺院、僧侣、道士和施主网络在材料中连接都城、州郡与边地补足资源与空间的视角。 它们不是由年份自动推出的因果，而是让前后事件之间的资源、道路、户籍或地方网络保持可见。 《隋书·高祖纪》；《陈书·后主纪》；《北史》；《资治通鉴》陈纪只能确认前一观察的基本年序，本纪与编年材料主要确证政权年序和精英行动；对基层执行、疆界和人口数量的沉默不得以叙事补足；《隋书》食货地理二志；墓志、寺院题记与江南城址为后一观察提供交叉线索，僧传、道教文本、造像记与遗址的成书或营造年代须分开，个案不能量化信仰或政策落实。


《隋书》《陈书》《北史》《资治通鉴》可证行动，志书、墓志、题记和城址只能检验局部。


\subsection{587—589年：材料所见的连续压力}
伐陈前后的转换要穿过长江渡运、州县户籍、寺院和家族网络；589年是最高政权转换日期，不是江南已同质化的证明。

沿着公元587年（年序桥接），隋与陈及突厥、岭南地方集团的\eventanchor{JH-DENS-587-A}{北周末隋初的制度接收与陈朝江南经营并行，淮河、长江和岭南水陆通道、户籍与军粮为再统一提供条件}指向北周末隋初的制度接收与陈朝江南经营并行，淮河、长江和岭南水陆通道、户籍与军粮为再统一提供条件，\eventanchor{JH-DENS-587-B}{隋与陈及突厥、岭南地方集团的关隘、渡口、军镇和水陆道路的可通行性继续约束使节、军队和物资的行动}则把隋与陈及突厥、岭南地方集团的关隘、渡口、军镇和水陆道路的可通行性继续约束使节、军队和物资的行动接入同一条年序。 它们不是由年份自动推出的因果，而是让前后事件之间的资源、道路、户籍或地方网络保持可见。 前者依据《隋书·高祖纪》；《陈书·后主纪》；《北史》；《资治通鉴》陈纪；本纪与编年材料主要确证政权年序和精英行动；对基层执行、疆界和人口数量的沉默不得以叙事补足。后者参照《隋书》食货地理二志；墓志、寺院题记与江南城址；纪传中的行军路线、兵数与控制范围常被简化或夸饰，地理材料只能提供有限交叉。

\eventanchor{JH-DENS-588-A}{北周末隋初的制度接收与陈朝江南经营并行，淮河、长江和岭南水陆通道、户籍与军粮为再统一提供条件}和\eventanchor{JH-DENS-588-B}{隋与陈及突厥、岭南地方集团的朝廷和州郡围绕户籍、田租、军粮与转运的安排进入材料}使公元588年（年序桥接）的隋与陈及突厥、岭南地方集团既保留北周末隋初的制度接收与陈朝江南经营并行，淮河、长江和岭南水陆通道、户籍与军粮为再统一提供条件，也不略去隋与陈及突厥、岭南地方集团的朝廷和州郡围绕户籍、田租、军粮与转运的安排进入材料。 这两项观察不宣称另有一场未载战争或政策，只把已有压力的时间位置和可见面向分开。 就可说范围而言，前一判断止于《隋书·高祖纪》；《陈书·后主纪》；《北史》；《资治通鉴》陈纪，本纪与编年材料主要确证政权年序和精英行动；对基层执行、疆界和人口数量的沉默不得以叙事补足；后一判断止于《隋书》食货地理二志；墓志、寺院题记与江南城址，志书、诏令与本纪可显示制度设计或战争需求，不能直接换算赋额、仓储或实际负担。

在公元589年（年序桥接），\eventanchor{JH-DENS-589-A}{北周末隋初的制度接收与陈朝江南经营并行，淮河、长江和岭南水陆通道、户籍与军粮为再统一提供条件}所见的北周末隋初的制度接收与陈朝江南经营并行，淮河、长江和岭南水陆通道、户籍与军粮为再统一提供条件与\eventanchor{JH-DENS-589-B}{隋与陈及突厥、岭南地方集团的流民、侨置户、军镇居民或地方家族的安置与编户问题构成社会背景}所见的隋与陈及突厥、岭南地方集团的流民、侨置户、军镇居民或地方家族的安置与编户问题构成社会背景并行发生在隋与陈及突厥、岭南地方集团的历史空间。 这两项观察不宣称另有一场未载战争或政策，只把已有压力的时间位置和可见面向分开。 《隋书·高祖纪》；《陈书·后主纪》；《北史》；《资治通鉴》陈纪使前一线索可追，本纪与编年材料主要确证政权年序和精英行动；对基层执行、疆界和人口数量的沉默不得以叙事补足；《隋书》食货地理二志；墓志、寺院题记与江南城址限制后一线索的说法，墓志、家族文本和侨置名目各自有选择性，不能据此统计迁徙规模或概括整体社会态度。


《隋书·高祖纪》《陈书·后主纪》《北史》《资治通鉴》可证陈亡，接收节奏仍须分开观察。
公元589年（材料补点）的材料将隋与陈、江南地方社会的\eventanchor{JH-DENS-589-X}{陈亡后的行政接收证据补点：隋军入建康后，最高政权转换与江南州县、寺院、家族和水路网络的实际接收需分开观察。}系于陈亡后的行政接收证据补点：隋军入建康后，最高政权转换与江南州县、寺院、家族和水路网络的实际接收需分开观察。，其影响由此进入后续年序。 它使原有的北方政权的军事胜利与南方地方社会的行政整合并非同一过程。转化为需要继续承担的提示589年是再统一的关键日期，而非所有制度和社会网络已经同质化的证明。 所据《隋书·高祖纪》；《陈书·后主纪》；墓志、寺院题记与江南城址材料留下可核的线索，同时提醒读者正史可确证陈亡与政权转换；地方户籍、寺院和家族网络的接收节奏需由不同材料分别检验。
